\chapter{Node.js}

Sappiamo ora che le pagine web possono essere generate al volo da programmi in
esecuzione sul server. E sappiamo già molto di JavaScript: finora, però, l'abbiamo
usato soltanto nell'ambiente del browser.

Nulla vieta di usarlo anche in un \textbf{ambiente server}. Esistono diversi
runtime che lo permettono --- \term{Node.js}, \term{Deno}, \term{Bun} --- e in
questo corso useremo il primo, di gran lunga il più diffuso.

\section{Che cos'è Node.js}

\dfn{Node.js}{
  \term{Node.js} è un \textbf{ambiente di esecuzione} (\textit{runtime})
  JavaScript, open source e multipiattaforma. Esegue il motore JavaScript
  \textbf{V8} di Google, lo stesso del progetto Chromium, e adotta un modello di
  I/O \textbf{basato sugli eventi, asincrono per impostazione predefinita e non
  bloccante}.}

È di gran lunga l'ambiente di esecuzione JavaScript più usato: nel sondaggio
\textit{State of JavaScript 2022} lo indicava il 93,6\% di chi ha risposto, più
ancora del browser.

\begin{figure}[H]
  \centering
  \begin{tikzpicture}[font=\Lato\scriptsize, x=0.085cm, y=0.5cm]
    % griglia e scala
    \foreach \p in {0,20,...,100} {
      \draw[neutral!22!white, line width=0.4pt] (\p,0.55) -- (\p,-6.55);
      \node[text=neutral!40!black, anchor=south] at (\p,0.55) {\p\%};
    }
    \foreach \nome/\val/\testo [count=\i from 0] in
        {Node.js/93.6/{93,6\%}, Browser/81.6/{81,6\%},
         Service Worker/17.8/{17,8\%}, Deno/11.2/{11,2\%}, Bun/4.3/{4,3\%},
         Hermes/3.0/{3,0\%}, ChakraCore/0.4/{0,4\%}} {
      \ifnum\i=0
        \fill[green-gradient-6, rounded corners=1pt] (0,-\i-0.3) rectangle (\val,-\i+0.3);
        \node[anchor=east, font=\Lato\scriptsize\bfseries] at (-1.5,-\i) {\nome};
        \node[anchor=west, font=\Lato\scriptsize\bfseries] at (\val+1,-\i) {\testo};
      \else
        \fill[primary!35!white, rounded corners=1pt] (0,-\i-0.3) rectangle (\val,-\i+0.3);
        \node[anchor=east] at (-1.5,-\i) {\nome};
        \node[anchor=west, text=neutral!30!black] at (\val+1,-\i) {\testo};
      \fi
    }
  \end{tikzpicture}
  \caption{Gli ambienti in cui si esegue JavaScript, secondo il sondaggio
           \textit{State of JavaScript 2022}. Si poteva indicarne più d'uno, per
           questo le percentuali non sommano a 100.}
  \label{fig:stateofjs}
\end{figure}

Per capire perché quel modello di I/O sia così importante, conviene partire da
come funzionano gli altri ambienti.

\section{Il modello multi-thread}

In molti ambienti e linguaggi, \textbf{una richiesta è gestita da un
thread o processo dedicato}, prelevato da un pool preallocato dall'application
server. Quel thread resta legato alla sua richiesta dall'inizio alla fine,
comprese le operazioni di I/O necessarie per servirla.

\begin{figure}[H]
  \centering
  \begin{tikzpicture}[font=\Lato\small,
      cl/.style={wtnode, minimum width=1.9cm, minimum height=0.7cm},
      rq/.style={rounded corners=2pt, draw=primary!65!black, fill=primary!14!white,
                 line width=0.7pt, minimum width=4.6cm, minimum height=0.72cm,
                 font=\Lato\scriptsize, text depth=0pt},
      th/.style={circle, draw=primary!65!black, fill=white, line width=0.7pt,
                 inner sep=0pt, minimum size=0.6cm, font=\Lato\scriptsize},
      io/.style={wtnodewarn, minimum width=2.5cm, minimum height=0.62cm,
                 font=\Lato\scriptsize},
      box/.style={rounded corners=5pt, line width=0.8pt},
      dbl/.style={{Stealth[length=5pt,width=4pt]}-{Stealth[length=5pt,width=4pt]},
                  line width=0.9pt, draw=primary!70!black}]

    % client
    \node[cl] (c1) at (0,1.1)  {Client 1};
    \node[cl] (c2) at (0,0.2)  {Client 2};
    \node[font=\Lato\small, text=neutral!45!black] at (0,-0.55) {\vdots};
    \node[cl] (cn) at (0,-1.4) {Client \textit{n}};

    % application server
    \draw[box, draw=green-gradient-6, fill=green-gradient-1!30!white]
      (2.0,2.2) rectangle (8.4,-3.55);
    \node[font=\Lato\scriptsize\bfseries, text=green-gradient-6!70!black,
          anchor=north] at (5.2,2.15) {Application server};
    \draw[box, dashed, draw=primary!55!black, fill=white]
      (2.3,1.65) rectangle (8.1,-1.95);
    \foreach \i/\y in {1/1.1, 2/0.2, n/-1.4} {
      \node[rq, anchor=west] at (2.55,\y)
        {\hspace*{-1.9cm}richiesta \textit{\i} gestita da};
      \node[th] at (7.35,\y) {T\textsubscript{\textit{\i}}};
      \draw[dbl] (c\i.east) -- (2.0,\y);
    }
    \node[font=\Lato\small, text=neutral!45!black] at (5.2,-0.55) {\vdots};
    \draw[box, dashed, draw=primary!55!black, fill=white]
      (2.3,-2.2) rectangle (8.1,-3.35);
    \node[font=\Lato\scriptsize\bfseries, text=primary!65!black, anchor=north]
      at (5.2,-2.22) {pool interno di thread / processi};
    \foreach \x/\t in {3.2/1, 4.1/2, 5.0/3, 7.2/\textit{m}} {
      \node[th, minimum size=0.48cm, font=\Lato\tiny] at (\x,-2.95) {T\textsubscript{\t}};
    }
    \node[text=neutral!45!black] at (6.1,-2.95) {$\cdots$};

    % operazioni di I/O
    \draw[box, draw=orange-gradient-6, fill=orange-gradient-1!25!white]
      (9.6,2.2) rectangle (12.6,-3.55);
    \node[font=\Lato\scriptsize\bfseries, text=orange-gradient-6!80!black,
          anchor=north] at (11.1,2.15) {Operazioni di I/O};
    \node[io, shape=cylinder, shape border rotate=90, aspect=0.25,
          minimum height=0.8cm, minimum width=2.3cm] at (11.1,1.0) {database};
    \node[io] at (11.1,-0.2) {filesystem};
    \node[io] at (11.1,-1.3) {API esterne};
    \node[io] at (11.1,-2.4) {dati di sistema};
    \draw[dbl] (8.4,-0.6) -- (9.6,-0.6);
  \end{tikzpicture}
  \caption{Il modello multi-thread: ogni richiesta occupa un thread del pool
           per tutta la sua durata, anche mentre aspetta le operazioni di I/O.}
  \label{fig:multithread}
\end{figure}

Il problema è che, per servire una richiesta, un'applicazione web deve spesso
compiere operazioni di I/O \textbf{lente e bloccanti}:

\begin{itemize}
  \item accedere a un database;
  \item accedere a un'API esterna (per esempio con \code{fetch()});
  \item leggere un file dal filesystem locale;
  \item leggere il corpo della richiesta.
\end{itemize}

Di conseguenza, i thread che gestiscono le richieste passano \textbf{molto
tempo ad aspettare} che queste operazioni si completino.

\begin{figure}[H]
  \centering
  \begin{tikzpicture}[font=\Lato\scriptsize,
      act/.style={fill=primary!30!white, draw=primary!60!black, line width=0.5pt,
                  minimum height=0.45cm, anchor=west, inner sep=0pt},
      wait/.style={fill=neutral!14!white, draw=neutral!45!white, line width=0.5pt,
                   minimum height=0.45cm, anchor=west, inner sep=0pt},
      io/.style={fill=orange-gradient-2!60!white, draw=orange-gradient-6,
                 line width=0.5pt, minimum height=0.38cm, anchor=west,
                 inner sep=0pt, font=\Lato\tiny, text=orange-gradient-6!70!black},
      ev/.style={rounded corners=2pt, fill=primary!75!black, text=white,
                 font=\Lato\scriptsize\bfseries, inner sep=3pt},
      fr/.style={-{Stealth[length=5pt,width=4pt]}, line width=0.9pt,
                 draw=primary!75!black}]

    % barra del thread
    \node[act,  minimum width=0.9cm] at (0,0) {};
    \node[wait, minimum width=2.0cm] at (0.9,0) {};
    \node[act,  minimum width=0.7cm] at (2.9,0) {};
    \node[wait, minimum width=2.4cm] at (3.6,0) {};
    \node[act,  minimum width=0.7cm] at (6.0,0) {};
    \node[wait, minimum width=1.8cm] at (6.7,0) {};
    \node[act,  minimum width=0.9cm] at (8.5,0) {};

    \node[font=\Lato\scriptsize, anchor=east, text=neutral!25!black]
      at (-0.2,0) {thread};

    % operazioni di I/O
    \node[io, minimum width=2.0cm] at (0.9,-0.75) {database};
    \node[io, minimum width=2.4cm] at (3.6,-0.75) {filesystem};
    \node[io, minimum width=1.8cm] at (6.7,-0.75) {database};
    \node[font=\Lato\scriptsize, anchor=east, text=neutral!25!black, align=right]
      at (-0.2,-0.75) {operazioni\\ di I/O};

    % etichette
    \node[font=\Lato\scriptsize, text=neutral!35!black] at (1.9,0.4) {attesa};
    \node[font=\Lato\scriptsize, text=neutral!35!black] at (4.8,0.4) {attesa};
    \node[font=\Lato\scriptsize, text=neutral!35!black] at (7.6,0.4) {attesa};

    \draw[-{Stealth[length=5pt]}, draw=neutral!45!black, line width=0.7pt]
      (0,-1.35) -- (9.6,-1.35);
    \node[font=\Lato\scriptsize, text=neutral!35!black, anchor=north]
      at (9.6,-1.42) {tempo};

    % la richiesta arriva, la risposta parte
    \node[ev] (rq) at (0.12,1.15) {richiesta};
    \draw[fr] (rq.south) -- (0.12,0.25);
    \node[ev] (rs) at (9.28,1.15) {risposta};
    \draw[fr] (9.28,0.25) -- (rs.south);
  \end{tikzpicture}
  \caption{Il thread che serve una richiesta: in azzurro il tempo in cui lavora
           davvero, in grigio quello in cui aspetta che finisca un'operazione di
           I/O \textbf{bloccante} (in arancione).}
  \label{fig:blocking-io}
\end{figure}

Questo modello ha due difetti:

\begin{itemize}
  \item è \textbf{inefficiente}: moltissimo tempo di thread è speso ad
        aspettare;
  \item ha \textbf{scalabilità limitata}: il pool di thread è preallocato a una
        dimensione fissa, e non si possono gestire più richieste concorrenti di
        quante ne consenta il pool.
\end{itemize}

\subsection{Una metafora: il ristorante}

Il docente propone un'analogia che vale la pena riportare per esteso, perché
rende il concetto immediato.

Un ristorante ha \textbf{10 tavoli}. La direzione assume \textbf{10 camerieri},
assegnandone uno a ciascun tavolo.

\begin{itemize}
  \item Quando arriva un cliente, viene fatto accomodare al primo tavolo libero
        e il suo cameriere gli porta il menu. Il cliente impiega dai 5 ai 10
        minuti per decidere: \textbf{in questo tempo il cameriere resta fermo}.
  \item Deciso l'ordine, il cameriere lo porta in cucina. Il cuoco impiega dai
        15 ai 30 minuti per preparare: \textbf{il cameriere resta fermo}.
  \item Quando il piatto è pronto, il cameriere lo porta al tavolo. Il cliente
        impiega dai 15 ai 30 minuti per mangiare: \textbf{il cameriere resta
        fermo}.
  \item Finito il pasto, il cameriere sparecchia e chiede il conto alla cassa,
        che impiega dai 2 ai 5 minuti: \textbf{il cameriere resta fermo}.
  \item Preparato il conto, il cameriere lo porta al cliente, che paga e se ne
        va. Solo allora quel cameriere può servire un nuovo cliente.
\end{itemize}

I camerieri sono i \textbf{thread}, i clienti sono le \textbf{richieste}. Al
massimo 10 richieste possono essere servite contemporaneamente, e le attese
(decidere, cucinare, mangiare, preparare il conto) sono le operazioni
\textbf{bloccanti}. Per servire una richiesta in più bisogna assumere un nuovo
cameriere.

\subsection{Il ristorante efficiente}

Proviamo un'organizzazione diversa: il ristorante assume \textbf{un solo
cameriere} per tutti e 10 i tavoli.

\begin{itemize}
  \item Quando arriva un cliente, il cameriere lo accoglie, lo fa accomodare e
        gli porta il menu. \textbf{Mentre il cliente sceglie, il cameriere è
        libero} di occuparsi degli altri.
  \item Quando il cliente ha deciso, il cameriere porta l'ordine in cucina.
        \textbf{Mentre il cuoco prepara, il cameriere torna in sala} ed è
        disponibile.
  \item Quando il piatto è pronto (e la campanella della cucina suona), il
        cameriere lo porta al tavolo. \textbf{Mentre il cliente mangia, il
        cameriere è di nuovo libero.}
\end{itemize}

Un solo cameriere serve dieci tavoli senza problemi, perché non passa mai il
tempo ad aspettare. La campanella della cucina è la chiave: è un \textbf{evento}
che segnala che un'operazione lenta è terminata.

\section{Il modello di Node.js}

Node.js è basato su una filosofia molto simile a quella del ristorante
efficiente.

\begin{itemize}
  \item Un programma Node.js gira in un \textbf{unico processo}: non serve
        creare nuovi thread per ogni richiesta.
  \item Il runtime fornisce nella libreria standard un insieme di
        \term{primitive di I/O asincrone} che impediscono al codice JavaScript
        di bloccarsi.
  \item La maggior parte delle librerie adotta il paradigma non bloccante
        ovunque possibile, con ampio uso del pattern delle \textbf{callback}.
  \item Le operazioni sincrone sono l'\textbf{eccezione}, non la norma.
\end{itemize}

\begin{figure}[H]
  \centering
  \begin{tikzpicture}[font=\Lato\small,
      cl/.style={wtnode, minimum width=1.7cm, minimum height=0.65cm,
                 font=\Lato\scriptsize},
      box/.style={rounded corners=5pt, line width=0.8pt},
      slot/.style={rounded corners=1.5pt, draw=primary!55!black,
                   fill=primary!12!white, line width=0.6pt, minimum width=1.55cm,
                   minimum height=0.36cm, font=\Lato\tiny, inner sep=0pt},
      vuoto/.style={slot, dashed, fill=white},
      io/.style={wtnodewarn, minimum width=2.2cm, minimum height=0.6cm,
                 font=\Lato\scriptsize},
      a/.style={-{Stealth[length=6pt,width=5pt]}, line width=1.4pt,
                draw=primary!70!black},
      dbl/.style={{Stealth[length=5pt,width=4pt]}-{Stealth[length=5pt,width=4pt]},
                  line width=0.9pt, draw=primary!70!black},
      giro/.style={-{Stealth[length=8pt,width=8pt]}, line width=3pt,
                   draw=primary!70!black},
      lab/.style={font=\Lato\scriptsize, text=neutral!20!black, align=center}]

    % client
    \node[cl] (c1) at (0,1.0)  {Client 1};
    \node[cl] (c2) at (0,0.15) {Client 2};
    \node[font=\Lato\small, text=neutral!45!black] at (0,-0.5) {\vdots};
    \node[cl] (cn) at (0,-1.3) {Client \textit{n}};

    % il runtime
    \draw[box, draw=green-gradient-6, fill=green-gradient-1!30!white]
      (1.4,2.25) rectangle (10.95,-2.45);
    \node[font=\Lato\small\bfseries, text=green-gradient-6!70!black, anchor=north]
      at (6.2,2.2) {Node.js};

    % coda degli eventi
    \draw[box, draw=primary!55!black, fill=white] (1.85,1.75) rectangle (3.85,-2.05);
    \node[lab, font=\Lato\scriptsize\bfseries, anchor=north] at (2.85,1.7)
      {coda degli\\eventi};
    \foreach \y in {0.45,0.0,-1.05,-1.5} \node[slot] at (2.85,\y) {};
    \node[font=\Lato\small, text=neutral!45!black] at (2.85,-0.45) {\vdots};
    \foreach \c/\y in {c1/1.0, c2/0.15, cn/-1.3} \draw[dbl] (\c.east) -- (1.85,\y);

    % l'event loop
    \def\cx{6.35} \def\cy{-0.15} \def\r{1.2}
    \draw[giro] (\cx,\cy) ++(160:\r) arc[start angle=160, end angle=22, radius=\r];
    \draw[giro] (\cx,\cy) ++(-20:\r) arc[start angle=-20, end angle=-158, radius=\r];
    \node[lab] at (\cx,\cy) {\textbf{Event loop}\\[-1pt]\tiny(un solo thread)};

    % worker thread
    \draw[box, draw=primary!55!black, fill=white] (8.7,1.75) rectangle (10.6,-2.05);
    \node[lab, font=\Lato\scriptsize\bfseries, anchor=north] at (9.65,1.7)
      {worker\\thread};
    \node[slot] at (9.65,0.45)  {filesystem};
    \node[slot] at (9.65,0.0)   {rete};
    \node[slot] at (9.65,-0.45) {database};
    \node[vuoto] at (9.65,-1.05) {};
    \node[vuoto] at (9.65,-1.5) {};

    % il giro: coda -> loop -> worker -> loop -> coda
    \draw[a] (3.95,0.65) -- (5.2,0.65);
    \draw[a] (7.5,0.65) -- node[lab, above=1pt]{I/O\\bloccante} (8.6,0.65);
    \draw[a] (8.6,-0.95) -- node[lab, below=1pt]{callback} (7.5,-0.95);
    \draw[a] (5.2,-0.95) -- (3.95,-0.95);

    % operazioni di I/O
    \draw[box, draw=orange-gradient-6, fill=orange-gradient-1!25!white]
      (11.55,2.25) rectangle (14.05,-2.45);
    \node[font=\Lato\scriptsize\bfseries, text=orange-gradient-6!80!black,
          anchor=north] at (12.8,2.2) {Operazioni di I/O};
    \node[io, shape=cylinder, shape border rotate=90, aspect=0.25,
          minimum height=0.75cm, minimum width=2.0cm] at (12.8,1.05) {database};
    \node[io] at (12.8,0.0)   {filesystem};
    \node[io] at (12.8,-0.85) {API esterne};
    \node[io] at (12.8,-1.7)  {dati di sistema};
    \draw[dbl] (10.6,-0.15) -- (11.55,-0.15);
  \end{tikzpicture}
  \caption{L'event loop di Node.js. Un solo thread esegue il codice JavaScript;
           l'I/O bloccante è affidato ai worker thread interni al runtime, e a
           operazione conclusa la sua callback torna all'event loop.}
  \label{fig:event-loop}
\end{figure}

Lo schema si legge come il ristorante efficiente. Le richieste dei client
diventano \textbf{eventi} in coda. L'\term{event loop} (il cameriere) preleva un
evento alla volta ed esegue il codice JavaScript che lo gestisce. Quando quel
codice avvia un'operazione di I/O, l'event loop non aspetta: la affida ai
\textbf{worker thread} del runtime (la cucina) e passa subito all'evento
successivo. A operazione conclusa, la \textbf{callback} che le era stata
associata (la campanella) viene rimessa in coda, e l'event loop la esegue appena
è libero.

\info{Un thread solo, ma non per tutto}{
  \guillemotleft Node.js è a thread singolo\guillemotright\ è una semplificazione utile ma
  imprecisa. È il \textbf{nostro} codice JavaScript a girare in un solo thread:
  le operazioni di I/O vere e proprie sono delegate a un pool di worker interni
  al runtime. La conseguenza pratica è importante: un calcolo lungo e
  \textbf{sincrono} nel nostro codice (un ciclo pesante, una cifratura,
  l'elaborazione di un'immagine) blocca l'event loop e quindi
  \textbf{tutte} le richieste, non solo la propria. Node.js eccelle nei carichi
  dominati dall'I/O, molto meno in quelli dominati dalla CPU.}

\section{Installare Node.js}

Ci sono tre modi:

\begin{itemize}
  \item usare gli installer disponibili sul sito ufficiale;
  \item usare il proprio gestore di pacchetti preferito;
  \item usare un \term{gestore di versioni} come \code{nvm}, \code{nvs} o
        \code{nvm4w}, comodo per sperimentare con versioni diverse e passare
        rapidamente dall'una all'altra.
\end{itemize}

Nel corso (A.A.\ 2026/27) si usa Node.js \textbf{22.14.0}; con qualunque
versione supportata non dovrebbero esserci problemi.

\subsection{Un gestore di versioni: nvm}

Con \code{nvm} installare una versione precisa è un solo comando. nvm scarica
l'archivio dal sito di Node.js, ne verifica l'integrità confrontando il
\textit{checksum} e passa subito a usarla. Insieme a Node.js arriva anche
\textbf{npm}, il gestore di pacchetti che vedremo fra poco, nella versione
abbinata.

\begin{console}[Terminale]
$ nvm install 22.14.0
Downloading and installing node v22.14.0...
Downloading https://nodejs.org/dist/v22.14.0/node-v22.14.0-darwin-arm64.tar.xz...
############################################################################ 100.0%
Computing checksum with sha256sum
Checksums matched!
Now using node v22.14.0 (npm v10.9.2)
\end{console}

Si possono tenere installate più versioni. \code{nvm list} le elenca: la
freccia indica quella in uso nel terminale corrente, \code{default} quella che
si attiva in ogni nuovo terminale. Per passare a un'altra versione si usa
\code{nvm use}, per esempio \code{nvm use 22.14.0}.

\begin{affianco}[0.36][c]
\begin{console}[Terminale]
$ nvm list
->     v22.14.0
       v26.8.2
default -> 26 (-> v26.8.2)
\end{console}
\risultato
\begin{immagine}
\begin{tikzpicture}[d/.style={mnota, anchor=west, text width=6.4cm}]
  \node[mcod, anchor=west, text=green-gradient-6!80!black] at (0,0) {->};
  \node[d] at (1.4,0)    {la versione in uso adesso, in questo terminale};
  \node[mcod, anchor=west, text=accent!70!black] at (0,-0.75) {default};
  \node[d] at (1.4,-0.75) {la versione predefinita (qui la 26), che ogni
                           nuovo terminale userà finché non la si cambia};
\end{tikzpicture}
\end{immagine}
\end{affianco}

\subsection{Verificare l'installazione}

Per controllare che tutto sia a posto si chiede a Node.js la sua versione. Se
invece si lancia \code{node} \textbf{senza argomenti}, si apre la \term{REPL}
(\textit{Read-Eval-Print Loop}): un interprete interattivo che legge
un'espressione JavaScript, la valuta e ne stampa il risultato, come la console
dei DevTools ma senza browser. Si esce con \code{.exit} (o premendo due volte
\code{CTRL+C}).

\begin{affianco}[0.5][t]
\begin{console}[Terminale]
$ node --version
v22.14.0
\end{console}
\risultato
\begin{console}[Terminale]
$ node
Welcome to Node.js v22.14.0.
Type ".help" for more information.
> "web".toUpperCase()
'WEB'
> .exit
\end{console}
\end{affianco}

\begin{esempio}{Hello World in Node.js}
\begin{affianco}[0.45][c]
\begin{lstlisting}[style=js, name=hello-world.js]
"use strict";
console.log("Hello Node.js!");
\end{lstlisting}
\risultato
\begin{console}[Terminale]
$ node hello-world.js
Hello Node.js!
\end{console}
\end{affianco}
\end{esempio}

\subsection{Differenze rispetto all'ambiente browser}

\begin{center}
\renewcommand{\arraystretch}{1.4}
\begin{tabular}{@{}>{\raggedright\arraybackslash}p{6.3cm}@{\hspace{1.0cm}}>{\raggedright\arraybackslash}p{6.3cm}@{}}
  \textbf{Nel browser} & \textbf{In Node.js} \\
  JavaScript serve soprattutto a manipolare il DOM: ci sono \code{window},
  \code{document}, le API della piattaforma web come i cookie e
  \code{localStorage}.
  & Questi oggetti e queste API \textbf{non esistono}. In compenso c'è la
    libreria standard con i suoi moduli per l'accesso al filesystem e alla
    rete: si possono leggere e scrivere file, e mettersi in ascolto su un socket
    per ricevere richieste. \\
\end{tabular}
\end{center}

\section{npm: il gestore di pacchetti}

Uno dei principali motori del successo di Node.js è \term{npm}. A settembre 2022
erano disponibili oltre 2,1 milioni di pacchetti: il più grande repository per
un singolo linguaggio al mondo. Esiste un pacchetto per (quasi) tutto.

npm viene installato insieme a Node.js. Per verificarlo:

\begin{affianco}[0.45][c]
\begin{lstlisting}[style=shell]
npm --version
\end{lstlisting}
\risultato
\begin{console}[Terminale]
10.9.2
\end{console}
\end{affianco}

\subsection{Creare un pacchetto}

Per creare un pacchetto npm nella directory corrente si usa \code{npm init}.
Poiché in questo corso lavoreremo con i \textbf{moduli ES}, useremo la variante:

\begin{lstlisting}[style=shell]
npm init es6
\end{lstlisting}

Il comando chiede alcune informazioni sul progetto (nome, versione, repository,
licenza, autore, punto di ingresso, comando di test) e al termine crea un file
\code{package.json} nella stessa directory, che contiene tutte le informazioni
sul pacchetto, comprese le dipendenze e i comandi per eseguirlo. Si può pensarlo
come la versione npm del \code{pom.xml} di Maven.

\begin{affianco}[0.6][c]
\begin{lstlisting}[style=json, name=package.json]
{
  "name": "hello-web-technologies",
  "version": "0.0.1",
  "description": "Web Technologies' first npm package",
  "main": "app.js",
  "type": "module",
  "scripts": {
    "test": "echo \"Error: no test specified\" && exit 1"
  },
  "keywords": ["greetings"],
  "author": "Luigi Libero Lucio Starace",
  "license": "MIT"
}
\end{lstlisting}
\risultato
\begin{immagine}[Le tre voci da conoscere.]
\begin{tikzpicture}[
    v/.style={rounded corners=3pt, draw=primary!60!black, fill=primary!8!white,
      line width=0.7pt, font=\FiraCodeNL\scriptsize, anchor=west, minimum width=1.9cm},
    d/.style={mnota, anchor=west, text width=3.4cm}]
  \node[v] (m) at (0,0)     {"main"};
  \node[d] at (2.1,0)       {il file che esporta ciò che il pacchetto offre};
  \node[v] (t) at (0,-1.0)  {"type"};
  \node[d] at (2.1,-1.0)    {\code{"module"} abilita \code{import} ed \code{export}};
  \node[v] (s) at (0,-2.0)  {"scripts"};
  \node[d] at (2.1,-2.0)    {comandi eseguibili con \code{npm run <nome>}};
\end{tikzpicture}
\end{immagine}
\end{affianco}

\begin{itemize}
  \item il \textbf{file principale} (\code{app.js} nell'esempio) viene eseguito
        quando il codice client importa il nostro pacchetto: di norma esporta le
        variabili e le funzioni pubbliche;
  \item la proprietà \code{scripts} permette di definire \textbf{attività} da
        riga di comando. Si eseguono con \code{npm run <nome>}; per le attività
        standard, come \code{start} e \code{test}, basta \code{npm <nome>}.
        L'attività \code{test} creata da \code{npm init} non fa granché: stampa
        un errore ed esce;
  \item \code{"type": "module"} è ciò che abilita i moduli ES, cioè la sintassi
        \code{import}/\code{export} vista nel capitolo 6.
\end{itemize}

Aggiungiamo un'attività \code{start}, che esegue il file principale:

\begin{affianco}[0.55][c]
\begin{lstlisting}[style=json, name=package.json]
"scripts": {
  "start": "node app.js",
  "test": "echo \"Error: no test specified\" && exit 1"
},
\end{lstlisting}
\begin{lstlisting}[style=js, name=app.js]
"use strict";
console.log("Hello Node.js!");
\end{lstlisting}
\risultato
\begin{console}[Terminale]
$ npm start

> hello-web-technologies@0.0.1 start
> node app.js

Hello Node.js!
\end{console}
\end{affianco}

npm stampa il nome e la versione del pacchetto, poi il comando associato
all'attività, e infine lo esegue.

\subsection{Installare dipendenze}

Per installare una dipendenza si usa \code{npm install <nomepacchetto>}.
Proviamo con \code{cowsay}:

\begin{lstlisting}[style=shell]
npm install cowsay
\end{lstlisting}

Succedono due cose:

\begin{enumerate}
  \item npm tiene traccia della nuova dipendenza, aggiungendo una sezione
        \code{dependencies} al file \code{package.json} (a sinistra);
  \item le dipendenze vengono scaricate nella directory \code{node\_modules}
        (a destra). Si noti che npm ha scaricato non solo il pacchetto
        richiesto, ma anche le sue dipendenze (ricorsivamente).
\end{enumerate}

\begin{affianco}[0.4][c]
\begin{lstlisting}[style=json]
"dependencies": {
  "cowsay": "^1.5.0"
}
\end{lstlisting}
\risultato
\begin{immagine}[Una dipendenza dichiarata, molte scaricate:\\
                 anche quelle di \code{cowsay}, ricorsivamente.]
\begin{tikzpicture}[
    dir/.style={mcod, anchor=west, inner sep=1.5pt},
    sel/.style={dir, fill=accent!14!white, draw=accent!55!black, rounded corners=2pt},
    ramo/.style={draw=neutral!50!black, line width=0.5pt}]
  \node[dir] (r) at (0,0) {node\_modules/};
  \foreach \n/\st [count=\i] in {cowsay/sel, get-stdin/dir, string-width/dir,
      strip-final-newline/dir, yargs/dir, \ldots/dir} {
    \node[\st] (n\i) at (0.5,-\i*0.48) {\n};
    \draw[ramo] (0.25,-\i*0.48+0.38) |- (n\i.west);
  }
  \node[mnota, anchor=west] at ([xshift=6pt]n1.east) {quella richiesta};
  \node[mnota, anchor=west, text=neutral!45!black] at ([xshift=6pt]n3.east) {le sue dipendenze};
\end{tikzpicture}
\end{immagine}
\end{affianco}

\begin{affianco}[0.48][c]
\begin{lstlisting}[style=js, name=app.js]
import cowsay from 'cowsay';

console.log(cowsay.think({
  text: "Hello Web Tech!"
}));
\end{lstlisting}
\risultato
\begin{console}[Terminale]
 _________________
( Hello Web Tech! )
 -----------------
        o   ^__^
         o  (oo)\_______
            (__)\       )\/\
                ||----w |
                ||     ||
\end{console}
\end{affianco}

\warning{node\_modules non va versionato}{
  Quando si distribuisce un pacchetto npm basta fornire il \textbf{proprio}
  codice e il descrittore \code{package.json}: la directory
  \code{node\_modules} \textbf{non} va aggiunta al controllo di versione. Chi
  riceve il progetto esegue \code{npm install} e npm ricostruisce l'intero
  albero delle dipendenze a partire dal \code{package.json}. È il motivo per cui
  \code{node\_modules/} è la prima riga di quasi ogni \code{.gitignore}.}

\section{Una prima applicazione web con Node.js}

\begin{esempio}{Un server HTTP in venti righe}\label{es:server-http}
\begin{affianco}[0.6][c]
\codicepiccolo
\begin{lstlisting}[style=js, name=app.js]
import http from 'http';

const PORT = 3000;

let server = http.createServer(function (request, response) {
  let course = "Web Technologies";

  response.writeHead(200, { "Content-Type": "text/html" });
  response.write(`<!DOCTYPE html>
    <html><body>
      <h1>Hello ${course}</h1>
      <p>La data corrente e' ${new Date().toDateString()}</p>
    </body></html>`);
  response.end();
}).listen(PORT);

console.log(`Web app in ascolto sulla porta ${PORT}`);
\end{lstlisting}
\risultato
\begin{browser}[localhost:3000]
  {\raggedright\Lora\bfseries\LARGE Hello Web Technologies\par}\vspace{6pt}
  La data corrente e' Mon Oct 05 2026
\end{browser}
\begin{immagine}[Ogni richiesta invoca la callback,\\ che scrive la risposta.]
\begin{tikzpicture}[
    att/.style={rounded corners=3pt, draw=primary!65!black, fill=primary!12!white,
      line width=0.7pt, font=\Lato\scriptsize\bfseries, minimum width=1.6cm,
      minimum height=0.5cm},
    req/.style={-{Stealth[length=5pt,width=4pt]}, line width=0.8pt, draw=primary!70!black},
    res/.style={-{Stealth[length=5pt,width=4pt]}, line width=0.8pt, draw=green-gradient-6},
    msg/.style={mcod, above}]
  \node[att] (b) at (0,0) {browser};
  \node[att] (s) at (4.3,0) {app.js};
  \draw[req] ([yshift=4pt]b.east) -- node[msg]{GET /} ([yshift=4pt]s.west);
  \draw[res] ([yshift=-4pt]s.west) -- node[mcod, below]{200 text/html} ([yshift=-4pt]b.east);
\end{tikzpicture}
\end{immagine}
\end{affianco}

Qui si vede bene la differenza rispetto al server statico del capitolo 2: non
c'è alcuna document root, non c'è alcun file HTML sul disco. Il documento viene
\textbf{costruito a ogni richiesta}, e la funzione passata a
\code{createServer} è una \textbf{callback}, invocata dall'event loop ogni volta
che arriva una richiesta.

Si riconoscono anche gli elementi del protocollo studiati nel primo capitolo:
il codice di stato \code{200}, l'header \code{Content-Type}, il corpo della
risposta.
\end{esempio}

\section{Strumenti di sviluppo}

\subsection{Debug in Visual Studio Code}

Il modo più semplice è abilitare la funzione \term{Auto Attach}, che collega un
debugger ai processi Node.js avviati dal terminale di VS Code. Per abilitarla si
apre la palette dei comandi (\code{CTRL+SHIFT+P} su Windows e Linux,
\code{CMD+SHIFT+P} su Mac), si cerca \guillemotleft Toggle Auto
Attach\guillemotright\ e si sceglie una modalità. Conviene la modalità
\textbf{Smart}, che si collega automaticamente solo agli script che non stanno
in \code{node\_modules}: il debugger si ferma nel nostro codice, non in quello
delle librerie.

\begin{figure}[H]
  \centering
  \begin{tikzpicture}[font=\Lato\scriptsize,
      voce/.style={anchor=west, text=vscfg, inner sep=0pt},
      desc/.style={anchor=west, text=vsctabdim, inner sep=0pt}]
    \def\w{12.2}
    \fill[vsctabbar, rounded corners=3pt] (0,0) rectangle (\w,-3.55);
    \node[text=vsctabdim] at (\w/2,-0.3) {Toggle auto attach on this machine};
    \draw[vscaccent, fill=vscbg, line width=0.6pt] (0.2,-0.55) rectangle (\w-0.2,-0.95);
    \fill[vsctabtext] (0.33,-0.64) rectangle (0.36,-0.86);
    \node[desc] at (0.25,-1.3) {Temporarily disable auto attach in this session};
    \foreach \v/\d [count=\i] in
        {Always/Auto attach to every Node.js process launched in the terminal,
         Smart/Auto attach when running scripts that aren't in a node\_modules folder,
         Only With Flag/Only auto attach when the `{-}{-}inspect' flag is given,
         Disabled/Auto attach is disabled and not shown in status bar} {
      \ifnum\i=2
        \fill[vscaccent!70!black] (0.1,-1.3-\i*0.45-0.2) rectangle (\w-0.1,-1.3-\i*0.45+0.2);
        \node[voce, text=vsctabtext, font=\Lato\scriptsize\bfseries]
          (v) at (0.25,-1.3-\i*0.45) {\v};
        \node[desc, text=vsctabtext] at ([xshift=5pt]v.east) {\d};
      \else
        \node[voce] (v) at (0.25,-1.3-\i*0.45) {\v};
        \node[desc] at ([xshift=5pt]v.east) {\d};
      \fi
    }
  \end{tikzpicture}
  \caption{Le modalità di Auto Attach, nella palette dei comandi di VS Code.}
  \label{fig:autoattach}
\end{figure}

Avviando l'applicazione dal terminale di VS Code, il debugger si collega da
solo. Lo si capisce dalle righe che Node.js stampa prima dell'output del
programma (\textit{Debugger attached}) e dal pannello \textit{Run and Debug},
che si attiva.

\begin{console}[Terminale di VS Code]
$ npm start

> hello-web-technologies@0.0.1 start
> node app.js

Debugger listening on ws://127.0.0.1:58896/06d4d583-e2b3-4a5f-baca-8635eb9dbedf
For help, see: https://nodejs.org/en/docs/inspector
Debugger attached.
Web app in ascolto sulla porta 3000
\end{console}

I \term{punti di interruzione} (\textit{breakpoint}) si aggiungono come di
consueto, facendo clic accanto al numero della riga: compare un pallino rosso.
Quando l'esecuzione raggiunge un breakpoint si \textbf{ferma}, e la riga su cui
è sospesa viene evidenziata. Il pannello \textit{Run and Debug} permette allora
di ispezionare le variabili in scope (\textit{Variables}), di osservare delle
espressioni scelte da noi (\textit{Watch}), di esaminare la pila delle chiamate
(\textit{Call Stack}) e molto altro.

% Il codice dell'editor e' un listato vero (con scheda e numeri di riga); il
% segno nel listato registra la posizione della riga 6, su cui si disegnano il
% breakpoint e l'evidenziazione.
\newcommand{\wtdbgsegno}{\tikz[remember picture, overlay]\coordinate (dbgriga);}
\newsavebox{\wtdbgeditor}
\begin{lrbox}{\wtdbgeditor}
\begin{minipage}{12.6cm}
\begin{lstlisting}[style=js, name=app.js, aboveskip=0pt, belowskip=0pt]
import http from 'http';

const PORT = 3000;

let server = http.createServer(function (request, response) {
  let course = "Web Technologies";

  (*@\wtdbgsegno@*)response.writeHead(200, { "Content-Type": "text/html" });
  response.write(`<!DOCTYPE html>
    <html><body>
      <h1>Hello ${course}</h1>
      <p>La data corrente e' ${new Date().toDateString()}</p>
    </body></html>`);
  response.end();
}).listen(PORT);

console.log(`Web app in ascolto sulla porta ${PORT}`);
\end{lstlisting}
\end{minipage}
\end{lrbox}

\begin{figure}[H]
  \centering
  \begin{tikzpicture}[remember picture, font=\Lato\scriptsize,
      sez/.style={anchor=west, inner sep=0pt, text=vscfg,
                  font=\Lato\fontsize{6}{7}\selectfont\bfseries},
      var/.style={anchor=west, inner sep=0pt, text=vscfg,
                  font=\FiraCodeNL\fontsize{6.2}{7}\selectfont}]
    \pgfmathsetlengthmacro{\h}{\ht\wtdbgeditor+\dp\wtdbgeditor}
    \def\sb{5.0}
    % barra laterale "Run and Debug" (prosegue sotto l'angolo dell'editor)
    \fill[vsctabbar] (\sb+0.3,-\h) -- (0,-\h) {[rounded corners=3pt] -- (0,0)
      -- (\sb+0.3,0)} -- cycle;
    \node[anchor=north west, inner sep=0pt] (ed) at (\sb,0) {\usebox{\wtdbgeditor}};
    \draw[vscbg, line width=0.6pt] (\sb,0) -- (\sb,-\h);

    \node[sez, text=vsctabdim] at (0.2,-0.3) {RUN AND DEBUG};
    \node[sez] at (0.2,-0.75) {\rotatebox[origin=c]{-90}{›}\ \ VARIABLES};
    \node[var] at (0.35,-1.1)  {\rotatebox[origin=c]{-90}{›}\ Local};
    \node[var] at (0.55,-1.42) {\ \ \textcolor{vscpurple}{course}: \textcolor{vscstring}{'Web Technologies'}};
    \node[var] at (0.55,-1.74) {› \textcolor{vscpurple}{request}: IncomingMessage \{…\}};
    \node[var] at (0.55,-2.06) {› \textcolor{vscpurple}{response}: ServerResponse \{…\}};
    \node[var] at (0.55,-2.38) {› \textcolor{vscpurple}{this}: Server};
    \node[var] at (0.35,-2.74) {\rotatebox[origin=c]{-90}{›}\ Module};
    \node[var] at (0.55,-3.06) {› \textcolor{vscpurple}{http}: \{createServer: ƒ, …\}};
    \node[var] at (0.55,-3.38) {\ \ \textcolor{vscpurple}{PORT}: \textcolor{vsccomment!60!vscfunc}{3000}};
    \node[var] at (0.55,-3.70) {› \textcolor{vscpurple}{server}: Server \{…\}};
    \node[sez] at (0.2,-4.16)  {\rotatebox[origin=c]{-90}{›}\ \ WATCH};
    \node[var] at (0.55,-4.48) {response.statusCode: \textcolor{vsccomment!60!vscfunc}{200}};
    \node[sez] at (0.2,-\h+0.75cm) {›\ \ CALL STACK};
    \node[sez] at (0.2,-\h+0.35cm) {›\ \ BREAKPOINTS};

    % riga in cui l'esecuzione e' sospesa, e il suo breakpoint
    \fill[yellow!80!orange, opacity=0.2]
      ([xshift=14pt, yshift=-2.8pt]ed.west |- dbgriga) rectangle
      ([yshift=7.4pt]ed.east |- dbgriga);
    \fill[red!85!black] ([xshift=3.4pt, yshift=2.4pt]ed.west |- dbgriga) circle (2.3pt);

    % barra delle azioni di debug
    \node[anchor=north east, fill=vsctabbar!70!vscfg, rounded corners=2pt,
          inner xsep=5pt, inner ysep=3pt]
      at ([xshift=-12pt, yshift=-2.5pt]ed.north east)
      {\wtdbgiconscura{maniglia}\hspace{6pt}\wtdbgiconscura{resume}\hspace{6pt}%
       \wtdbgiconscura{stepover}\hspace{6pt}\wtdbgiconscura{stepinto}\hspace{6pt}%
       \wtdbgiconscura{stepout}\hspace{6pt}\wtdbgiconscura{restart}\hspace{6pt}%
       \wtdbgiconscura{disconnect}};
  \end{tikzpicture}
  \caption{VS Code fermo su un breakpoint alla riga 8 dell'\code{app.js}
           dell'esempio~\ref{es:server-http}. A sinistra le variabili in scope (le
           locali della callback e quelle del modulo) e un'espressione
           osservata; in alto, la barra delle azioni di debug.}
  \label{fig:debug}
\end{figure}

Con le \textbf{azioni di debug}, i pulsanti della barra in alto, si controlla
il flusso dell'esecuzione un passo alla volta:

\begin{center}
\renewcommand{\arraystretch}{1.5}
\begin{tabular}{@{}c@{\hspace{8pt}}l@{\hspace{0.8cm}}>{\raggedright\arraybackslash}p{9.6cm}@{}}
  & \textbf{Azione} & \textbf{Descrizione} \\
  \wtdbgicon{resume}     & \textit{Resume}     & riprende il normale flusso di
                           esecuzione, fino al prossimo breakpoint. \\
  \wtdbgicon{stepover}   & \textit{Step Over}  & esegue l'istruzione successiva
                           come un blocco unico, senza entrare nei suoi passi
                           interni (se è una chiamata, la funzione viene eseguita
                           per intero). \\
  \wtdbgicon{stepinto}   & \textit{Step Into}  & entra nell'istruzione
                           successiva per seguirne l'esecuzione riga per riga. \\
  \wtdbgicon{stepout}    & \textit{Step Out}   & essendo dentro una funzione, ne
                           completa le righe restanti come se fossero un'unica
                           istruzione e torna al contesto di esecuzione
                           precedente. \\
  \wtdbgicon{restart}    & \textit{Restart}    & termina l'esecuzione e
                           ricomincia il debug con la configurazione corrente
                           (\textbf{non} funziona con Auto Attach). \\
  \wtdbgicon{disconnect} & \textit{Disconnect} & termina la sessione di debug. \\
\end{tabular}
\end{center}

\info{Meglio del \code{console.log}}{
  È molto comune, all'inizio, cercare di capire un bug disseminando il codice di
  \code{console.log}. Funziona, ma è lento e invasivo. Mezz'ora spesa a imparare
  il debugger si ripaga molte volte: si vedono \textbf{tutte} le variabili in
  scope in un colpo solo, senza modificare il codice e senza dover indovinare in
  anticipo quali stampare.}

\subsection{Ricaricamento automatico}

Ogni volta che modifichiamo il codice dobbiamo riavviare l'intero server
(\code{npm start}, oppure \code{node app.js}). Per rendere lo sviluppo più
piacevole si usano utilità che monitorano il codice e riavviano il server di
sviluppo quando serve: è il \term{live reloading}.

Una di queste è \term{nodemon}:

\begin{lstlisting}[style=shell]
npm install nodemon
\end{lstlisting}

Basta poi sostituire \code{node app.js} con \code{nodemon app.js}, per
esempio nell'attività \code{start}:

\begin{affianco}[0.36][c]
\begin{lstlisting}[style=json, name=package.json]
"scripts": {
  "start": "nodemon app.js"
},
\end{lstlisting}
\risultato
\begin{console}[Terminale di VS Code]
$ npm start

> hello-web-technologies@0.0.1 start
> nodemon app.js

[nodemon] 3.0.1
[nodemon] to restart at any time, enter `rs`
[nodemon] watching path(s): *.*
[nodemon] watching extensions: js,mjs,cjs,json
[nodemon] starting `node app.js`
Debugger attached.
Web app in ascolto sulla porta 3000
[nodemon] restarting due to changes...
[nodemon] starting `node app.js`
Debugger attached.
Web app in ascolto sulla porta 3000
\end{console}
\end{affianco}

All'avvio nodemon dice che cosa sorveglia: tutti i file della directory
(\code{*.*}) con estensione \code{js}, \code{mjs}, \code{cjs} o \code{json}.
Appena uno di questi cambia, termina il processo e lo riavvia
(\textit{restarting due to changes}); per un riavvio manuale basta scrivere
\code{rs} nel terminale e premere invio. Con Auto Attach attivo, il debugger
di VS Code si ricollega da solo a ogni riavvio.

\section{Riferimenti}

\begin{itemize}
  \item \textbf{Getting started} --- Node.js Docs. \\
        \urlx{https://nodejs.org/en/learn/} \\
        Parti rilevanti: \textit{Introduction to Node.js}, \textit{How to
        install Node.js}, \textit{Differences between Node.js and the Browser},
        \textit{The V8 JavaScript Engine}, \textit{An introduction to the npm
        package manager}, \textit{ECMAScript 2015 (ES6) and beyond}.
  \item \textbf{Eloquent JavaScript} (3\textsuperscript{a} edizione), Marijn
        Haverbeke, capitolo 20. \\
        \urlx{https://eloquentjavascript.net/}
  \item \textbf{Mixu's Node book}, Mikito Takada. \\
        \urlx{http://book.mixu.net/node/single.html} \\
        Parti rilevanti: capitolo 2 (\textit{What is Node.js?}), capitolo 8
        sezione 8.2 (\textit{NPM}).
  \item \textbf{Nodemon} --- sito ufficiale. \\
        \urlx{https://nodemon.io/}
  \item \textbf{Node.js debugging in VS Code}. \\
        \urlx{https://code.visualstudio.com/docs/nodejs/nodejs-debugging}
\end{itemize}
